\documentclass[../../main.tex]{subfiles}

\begin{document}

\section{内积和正交性}\label{section:ch16-inner-product}

\begin{definition}[内积空间]\label{definition:ch16-inner-product}
令 $H$ 为线性空间. 函数 $\langle\cdot,\cdot\rangle:H\times H\to\mathbf{R}$ 称为 $H$ 上的内积, 若对所有 $x_{1},x_{2},x$ 与 $y\in H$ 以及实数 $\alpha$ 和 $\beta$,
\begin{enumerate}[(i)]
\item $\langle\alpha x_{1}+\beta x_{2},\ y\rangle=\alpha\langle x_{1},\ y\rangle+\beta\langle x_{2},\ y\rangle$.
\item $\langle x,\ y\rangle=\langle y,\ x\rangle$.
\item 若 $x\neq 0$, $\langle x,\ x\rangle>0$.
\end{enumerate}
具有内积的线性空间 $H$ 称为内积空间.
\end{definition}

性质(ii)称为对称性. 从(i)和(ii)得出 $\langle x,\ \alpha y_{1}+\beta y_{2}\rangle=\alpha\langle x,\ y_{1}\rangle+\beta\langle x,\ y_{2}\rangle$: 该性质与(i)一起称为双线性.

在无穷维线性空间中我们立即想起两个内积空间的例子. 对于两个序列 $x=\{x_{k}\}$ 与 $y=\{y_{k}\}\in\ell^{2}$, $\ell^{2}$ 内积 $\langle x,\ y\rangle$ 定义为
\begin{align*}
\langle x,y\rangle=\sum_{k=1}^{\infty}x_{k}y_{k}.
\end{align*}
对 $E$ 是实数的可测集以及两个函数 $f$ 和 $g\in L^{2}(E)$, $L^{2}$ 内积 $\langle f,\ g\rangle$ 定义为
\begin{align*}
\langle f,g\rangle=\int_{E}f\cdot g.
\end{align*}
其中积分是关于 Lebesgue 测度的积分. 从我们在第 7 章得到的作为 H\"older 不等式的特殊情形的 Cauchy-Schwartz 不等式, 我们推出内积是恰当定义的. Cauchy-Schwartz 不等式对一般的内积空间成立.

\begin{theorem}[Cauchy-Schwartz 不等式]\label{theorem:ch16-cauchy-schwartz}
对于内积空间 $H$ 中的任何两个向量 $u,\ v$,
\begin{align*}
|\langle u,v\rangle|\leqslant\|u\|\cdot\|v\|.
\end{align*}
\end{theorem}

为证明它, 观察到:
\begin{align*}
\text{对所有 }t\in\mathbf{R},\quad 0\leqslant\|u+tv\|^{2}=\|u\|^{2}+2t\langle u,v\rangle+t^{2}\|v\|^{2}.
\end{align*}
右边定义的 $t$ 的二次多项式没有两个不同的实根, 因此它的判别式不是正的, 即 Cauchy-Schwartz 不等式成立.

\begin{proposition}\label{proposition:ch16-norm}
对内积空间 $H$ 的向量 $h$, 定义
\begin{align*}
\|h\|=\sqrt{\langle h,h\rangle}.
\end{align*}
则 $\|\cdot\|$ 是 $H$ 上的范数, 它称为由内积 $\langle\cdot,\ \cdot\rangle$ 诱导的范数.
\end{proposition}

\begin{proof}
对 $\|\cdot\|$ 来说, 仅有范数的三角不等式这一性质不是显然的. 然而它是 Cauchy-Schwartz 不等式的推论, 对两个向量 $u,\ v\in H$,
\begin{align*}
\|u+v\|^{2}=\langle u+v,u+v\rangle=\|u\|^{2}+2\langle u,v\rangle+\|v\|^{2}\leqslant\|u\|^{2}+2\|u\|\|v\|+\|v\|^{2}=(\|u\|+\|v\|)^{2}.
\end{align*}
\end{proof}

以下等式刻画了由内积诱导的范数, 见习题\cref{exercise:ch16-07}.

\begin{theorem}[平行四边形等式]\label{theorem:ch16-parallelogram}
对于内积空间 $H$ 的任何两个向量 $u,\ v$,
\begin{align*}
\|u-v\|^{2}+\|u+v\|^{2}=2\|u\|^{2}+2\|v\|^{2}.
\end{align*}
\end{theorem}

为证明该等式仅须把以下两个等式相加:
\begin{align*}
\|u-v\|^{2}=\|u\|^{2}-2\langle u,v\rangle+\|v\|^{2},\\
\|u+v\|^{2}=\|u\|^{2}+2\langle u,v\rangle+\|v\|^{2}.
\end{align*}

\begin{definition}[Hilbert 空间]\label{definition:ch16-hilbert}
内积空间 $H$ 称为 Hilbert 空间, 若它是 Banach 空间且范数由内积诱导.
\end{definition}

Riesz-Fischer 定理告诉我们对于实数的可测集 $E$, $L^{2}(E)$ 是 Hilbert 空间, 作为一个推论, $\ell^{2}$ 也如此.

\begin{proposition}\label{proposition:ch16-nearest-point}
令 $K$ 为 Hilbert 空间 $H$ 的非空闭凸子集而 $h_{0}$ 属于 $H\sim K$. 则恰好存在一个向量 $h_{*}\in K$, 它在
\begin{align*}
\|h_{0}-h_{*}\|=\operatorname{dist}(h_{0},K)=\inf_{h\in K}\|h_{0}-h\|
\end{align*}
的意义下离 $h_{0}$ 最近.
\end{proposition}

\begin{proof}
通过用 $K$ 代替 $K-h_{0}$, 我们可以假设 $h_{0}=0$. 令 $\{h_{n}\}$ 为 $K$ 中的序列使得
\begin{align}
\lim_{n\to\infty}\|h_{n}\|=\inf_{h\in K}\|h\|.\label{align:ch16-nearest-1}
\end{align}
我们从平行四边形等式与 $K$ 的凸性推出, 对每个 $m$ 和 $n$,
\begin{align}
\|h_{n}\|^{2}+\|h_{m}\|^{2}=2\left\|\frac{h_{n}+h_{m}}{2}\right\|^{2}+2\left\|\frac{h_{n}-h_{m}}{2}\right\|^{2}\geqslant 2\cdot\inf_{h\in K}\|h\|^{2}+2\cdot\left\|\frac{h_{n}-h_{m}}{2}\right\|^{2}.\label{align:ch16-nearest-2}
\end{align}
从\eqref{align:ch16-nearest-1}和\eqref{align:ch16-nearest-2}我们推出 $\{h_{n}\}$ 是 Cauchy 列. 由于 $H$ 是完备的且 $K$ 是闭的, $\{h_{n}\}$ 强收敛到 $h_{*}\in K$. 根据范数的连续性, $\|h_{*}\|=\inf_{h\in K}\|h\|$. $K$ 中离原点最近的点是唯一的. 的确, 若 $h_{\star}$ 是 $K$ 中离原点最近的其他点, 则我们在不等式\eqref{align:ch16-nearest-2}中用 $h_{\star}$ 代替 $h_{n}$, 用 $h_{*}$ 代替 $h_{m}$, 则有
\begin{align*}
0=\|h^{*}\|^{2}+\|h_{\star}\|^{2}-2\cdot\inf_{h\in K}\|h\|^{2}\geqslant 2\cdot\left\|\frac{h^{*}-h_{\star}}{2}\right\|^{2}.
\end{align*}
因此 $h^{*}=h_{\star}$.
\end{proof}

\begin{definition}[正交]\label{definition:ch16-orthogonal}
内积空间 $H$ 中的两个向量 $u,\ v$ 称为是正交的, 若 $\langle u,\ v\rangle=0$. 向量 $u$ 称为正交于 $H$ 的子集 $S$, 若它正交于 $S$ 中的每个向量. 我们用 $S^{\perp}$ 表示 $H$ 中与 $S$ 正交的向量的全体.
\end{definition}

我们把从 Cauchy-Schwartz 不等式推出若 $S$ 是内积空间 $H$ 的子集, 则 $S^{\perp}$ 是 $H$ 的闭子集留作练习. 以下定理是基本的.

\begin{theorem}[正交直和分解]\label{theorem:ch16-orthogonal-decomposition}
令 $V$ 为 Hilbert 空间 $H$ 的闭子空间. 则 $H$ 有正交直和分解:
\begin{align}
H=V\oplus V^{\perp}.\label{align:ch16-decomposition}
\end{align}
\end{theorem}

\begin{proof}
令 $h_{0}$ 属于 $H\sim V$. 前一个命题告诉我们存在唯一向量 $h^{*}\in V$ 离 $h_{0}$ 最近. 令 $h$ 为 $V$ 中的任何向量. 对于实数 $t$, 由于 $V$ 是线性空间, 向量 $h^{*}-th$ 属于 $V$, 因此,
\begin{align*}
\langle h_{0}-h^{*},h_{0}-h^{*}\rangle=\|h_{0}-h\|^{2}\leqslant\|h_{0}-(h^{*}-th)\|^{2}=\langle h_{0}-h^{*},h_{0}-h^{*}\rangle+2t\cdot\langle h_{0}-h^{*},h\rangle+t^{2}\langle h,h\rangle.
\end{align*}
因此,
\begin{align*}
\text{对所有 }t\in\mathbf{R},\quad 0\leqslant 2t\cdot\langle h_{0}-h^{*},h\rangle+t^{2}\langle h,h\rangle.
\end{align*}
从而 $\langle h_{0}-h^{*},\ h\rangle=0$. 因此向量 $h_{0}-h^{*}$ 正交于 $V$. 观察到 $h_{0}=h^{*}+[h_{0}-h^{*}]$. 我们得出 $H=V+V^{\perp}$, 且由于 $V\cap V^{\perp}=\{0\}$, $H=V\oplus V^{\perp}$. 我们把以下推论的证明留作练习.
\end{proof}

\begin{corollary}\label{corollary:ch16-complete-span}
令 $S$ 为 Hilbert 空间 $H$ 的子集. 则 $S$ 的闭线性扩张是整个 $H$ 当且仅当 $S^{\perp}=\{0\}$.
\end{corollary}

根据\eqref{align:ch16-decomposition}, 对 $H$ 的闭子空间 $V$, 我们称 $V^{\perp}$ 为 $V$ 在 $H$ 中的正交补, 而称\eqref{align:ch16-decomposition}为 $H$ 的正交分解. 算子 $P\in\mathcal{L}(H)$ 是 $H$ 沿着 $V^{\perp}$ 到 $V$ 的投影, 称为正交投影.

\begin{proposition}\label{proposition:ch16-orth-projection}
令 $P$ 为 Hilbert 空间 $H$ 映上 $H$ 的非平凡闭子空间 $V$ 的正交投影. 则 $\|P\|=1$ 且
\begin{align}
\text{对所有 }u,v\in H,\quad\langle Pu,v\rangle=\langle u,Pv\rangle.\label{align:ch16-projection}
\end{align}
\end{proposition}

\begin{proof}
令向量 $u$ 属于 $H$. 则
\begin{align*}
\|u\|^{2}=\langle P(u)+(\operatorname{Id}-P)(u),P(u)+(\operatorname{Id}-P)(u)\rangle=\|P(u)\|^{2}+\|(\operatorname{Id}-P)(u)\|^{2}\geqslant\|P(u)\|^{2}.
\end{align*}
因此 $\|P(u)\|\leqslant\|u\|$. 我们因此有 $\|P\|\leqslant 1$. 由于对 $V$ 中的每个非零向量 $v$, $P(v)=v$, 我们得出 $\|P\|=1$. 若向量 $v$ 也属于 $H$, 则
\begin{align*}
\langle P(u),(\operatorname{Id}-P)(v)\rangle=\langle(\operatorname{Id}-P)(u),P(v)\rangle=0.
\end{align*}
因此,
\begin{align*}
\langle P(u),v\rangle=\langle P(u),P(v)\rangle=\langle u,P(v)\rangle.
\end{align*}
\end{proof}

关于 Banach 空间的许多结果的证明在 Hilbert 空间的特殊情形下简单得多, 见习题\cref{exercise:ch16-11}--\cref{exercise:ch16-15}.

\medskip\par\noindent\textbf{注. }Banach 空间 $X$ 称为是可补的, 若 $X$ 的每个闭子空间有闭的线性补. Banach 空间 $X$ 称为是可 Hilbert 的, 若存在 $X$ 上的内积诱导的范数等价于给定的范数. 我们从定理\cref{theorem:ch16-orthogonal-decomposition}推出一个可 Hilbert 的 Banach 空间是可补的. Joram Lindenstrauss 与 Lior Tzafriri 的一个非凡的定理断言逆定理成立: 若 Banach 空间是可补的, 则它是可 Hilbert 的.

\begin{exercise}\label{exercise:ch16-01}
令 $[a,\ b]$ 为实数的闭有界区间. 证明 $L^{2}[a,\ b]$ 内积也是 $C[a,\ b]$ 上的内积. 具有 $L^{2}[a,\ b]$ 内积的空间 $C[a,\ b]$ 是 Hilbert 空间吗?
\end{exercise}

\begin{exercise}\label{exercise:ch16-02}
证明 $C[a,\ b]$ 上的最大值范数不是由内积诱导的. $\ell^{1}$ 上的通常范数也不是由内积诱导的.
\end{exercise}

\begin{exercise}\label{exercise:ch16-03}
令 $H_{1}$ 和 $H_{2}$ 为 Hilbert 空间. 证明具有使得 $H_{1}\times\{0\}=[\{0\}\times H_{2}]^{\perp}$ 的内积的笛卡儿积 $H_{1}\times H_{1}$ 也是一个 Hilbert 空间.
\end{exercise}

\begin{exercise}\label{exercise:ch16-04}
证明: 若 $S$ 是内积空间 $H$ 的子集, 则 $S^{\perp}$ 是 $H$ 的闭子空间.
\end{exercise}

\begin{exercise}\label{exercise:ch16-05}
令 $S$ 为 $H$ 的子集. 证明 $S=(S^{\perp})^{\perp}$ 当且仅当 $S$ 是 $H$ 的闭子空间.
\end{exercise}

\begin{exercise}\label{exercise:ch16-06}
(极化等式)证明对任何两个向量 $u,\ v\in H$,
\begin{align*}
\langle u,v\rangle=\frac{1}{4}\left[\|u+v\|^{2}-\|u-v\|^{2}\right].
\end{align*}
\end{exercise}

\begin{exercise}\label{exercise:ch16-07}
(Jordan-von Neumann)令 $X$ 为赋予 $\|\cdot\|$ 的线性空间. 用极化等式证明范数 $\|\cdot\|$ 由内积诱导当且仅当平行四边形等式成立.
\end{exercise}

\begin{exercise}\label{exercise:ch16-08}
令 $V$ 为 $H$ 的闭子空间而 $P$ 是 $H$ 映上 $V$ 的投影. 证明 $P$ 是 $H$ 映上 $V$ 的正交投影当且仅当\eqref{align:ch16-projection}成立.
\end{exercise}

\begin{exercise}\label{exercise:ch16-09}
令 $T$ 属于 $\mathcal{L}(H)$. 证明 $T$ 是等距当且仅当
\begin{align*}
\text{对所有 }u,v\in H,\quad\langle T(u),T(v)\rangle=\langle u,v\rangle.
\end{align*}
\end{exercise}

\begin{exercise}\label{exercise:ch16-10}
令 $V$ 为 $H$ 的有限维子空间, 而 $\varphi_{1},\ \cdots,\ \varphi_{n}$ 是由单位向量组成的基. 基的每对元素是正交的. 证明 $H$ 的映上 $V$ 的正交投影 $P$ 由
\begin{align*}
P(h)=\sum_{k=1}^{n}\langle h,\varphi_{k}\rangle\varphi_{k},\quad h\in V
\end{align*}
给出.
\end{exercise}

\begin{exercise}\label{exercise:ch16-11}
对 $H$ 中的向量 $h$, 证明函数 $u\mapsto\langle h,\ u\rangle$ 属于 $H^{*}$.
\end{exercise}

\begin{exercise}\label{exercise:ch16-12}
对任何向量 $h\in H$, 证明存在有界线性泛函 $\psi\in H^{*}$ 使得
\begin{align*}
\|\psi\|=1\text{ 且 }\psi(h)=\|h\|.
\end{align*}
\end{exercise}

\begin{exercise}\label{exercise:ch16-13}
令 $V$ 为 $H$ 的闭子空间而 $P$ 是 $H$ 映上 $V$ 的正交投影. 对任何赋范线性空间 $X$ 与 $T\in\mathcal{L}(V,\ X)$, 证明 $T\circ P$ 属于 $\mathcal{L}(H,\ X)$, 且它是 $T:V\to X$ 的延拓, 使得 $\|T\circ P\|=\|T\|$.
\end{exercise}

\begin{exercise}\label{exercise:ch16-14}
直接用命题\cref{proposition:ch16-nearest-point}证明关于 $H$ 的超平面分离定理, 这里 $H$ 考虑为具有强拓扑的局部凸拓扑向量空间.
\end{exercise}

\begin{exercise}\label{exercise:ch16-15}
用命题\cref{proposition:ch16-nearest-point}证明 Hilbert 空间的 Krein-Milman 引理.
\end{exercise}

\end{document}